\documentclass{beamer}

\pdfmapfile{+sansmathaccent.map}


\mode<presentation>
{
	\usetheme{Warsaw} % or try Darmstadt, Madrid, Warsaw, Rochester, CambridgeUS, ...
	\usecolortheme{seahorse} % or try seahorse, beaver, crane, wolverine, ...
	\usefonttheme{serif}  % or try serif, structurebold, ...
	\setbeamertemplate{navigation symbols}{}
	\setbeamertemplate{caption}[numbered]
} 


%%%%%%%%%%%%%%%%%%%%%%%%%%%%
% itemize settings


%%%%%%%%%%%%%%%%%%%%%%%%%%%%
% itemize settings

\definecolor{myhotpink}{RGB}{255, 80, 200}
\definecolor{mywarmpink}{RGB}{255, 60, 160}
\definecolor{mylightpink}{RGB}{255, 80, 200}
\definecolor{mypink}{RGB}{255, 30, 80}
\definecolor{mydarkpink}{RGB}{155, 25, 60}

\definecolor{mypaleblue}{RGB}{240, 240, 255}
\definecolor{mylightblue}{RGB}{120, 150, 255}
\definecolor{myblue}{RGB}{90, 90, 255}
\definecolor{mygblue}{RGB}{70, 110, 240}
\definecolor{mycourseblue}{RGB}{120, 140, 255}
\definecolor{mycourselightblue}{RGB}{245, 250, 255}
\definecolor{mydarkblue}{RGB}{0, 0, 180}
\definecolor{myblackblue}{RGB}{40, 40, 120}

\definecolor{mygreen}{RGB}{0, 200, 0}
\definecolor{mygreen2}{RGB}{245, 255, 230}
\definecolor{mydarkgreen}{RGB}{0, 120, 0}
\definecolor{mydarkgreen}{RGB}{0, 120, 0}


\definecolor{mygray}{gray}{0.8}
\definecolor{mydarkgray}{RGB}{80, 80, 160}
\definecolor{mylightgray}{RGB}{160, 160, 160}

\definecolor{mydarkred}{RGB}{160, 30, 30}
\definecolor{mylightred}{RGB}{255, 150, 150}
\definecolor{myred}{RGB}{200, 110, 110}
\definecolor{myblackred}{RGB}{120, 40, 40}

\definecolor{mypink}{RGB}{255, 30, 80}
\definecolor{myhotpink}{RGB}{255, 80, 200}
\definecolor{mywarmpink}{RGB}{255, 60, 160}
\definecolor{mylightpink}{RGB}{255, 80, 200}
\definecolor{mydarkpink}{RGB}{155, 25, 60}
\definecolor{mywhitepink}{RGB}{255, 240, 240}

\definecolor{mydarkcolor}{RGB}{60, 25, 155}
\definecolor{mylightcolor}{RGB}{130, 180, 250}



\setbeamertemplate{itemize items}[default]

\setbeamertemplate{itemize item}{\color{myblackblue}$\blacksquare$}
\setbeamertemplate{itemize subitem}{\color{mygblue}$\blacktriangleright$}
\setbeamertemplate{itemize subsubitem}{\color{mygray}$\blacksquare$}

\setbeamercolor{palette quaternary}{fg=white,bg=mycourseblue}
\setbeamercolor{titlelike}{parent=palette quaternary}

\setbeamercolor{palette quaternary2}{fg=black,bg=mycourselightblue}
\setbeamercolor{frametitle}{parent=palette quaternary2}

\setbeamerfont{frametitle}{size=\Large,series=\scshape}
\setbeamerfont{framesubtitle}{size=\normalsize,series=\upshape}





%%%%%%%%%%%%%%%%%%%%%%%%%%%%
% block settings

\setbeamercolor{block title}{bg=red!30,fg=black}

\setbeamercolor*{block title example}{bg=mygreen!40!white,fg=black}

\setbeamercolor*{block body example}{fg= black, bg= mygreen2}


%%%%%%%%%%%%%%%%%%%%%%%%%%%%
% URL settings
\hypersetup{
	colorlinks=true,
	linkcolor=blue,
	filecolor=blue,      
	urlcolor=blue,
}

%%%%%%%%%%%%%%%%%%%%%%%%%%

\renewcommand{\familydefault}{\rmdefault}


\newcommand{\mycolor}[1] {
\setbeamercolor{normal text}{fg={#1}} % Recolors standard body text
\setbeamercolor{structure}{fg={#1}}   % Recolors bullets and structure items
\usebeamercolor[fg]{normal text}      % Applies the normal text color immediately
}


\usepackage{amsmath}
\usepackage{mathtools}

\usepackage{cancel}
\usepackage{centernot}

\usepackage{changepage}
\usepackage{subcaption}

\usepackage{qrcode}

\newcommand{\bo}[1] {\mathbf{#1}}
\newcommand{\R}{\mathbb{R}} 
\newcommand{\T}{^\top}     



\newcommand{\mydate}{Fall 2026}

\newcommand{\mygit}{\textcolor{blue}{\href{https://github.com/SergeiSa/Advanced-Control-2026}{github.com/SergeiSa/Advanced-Control-2026}}}

\newcommand{\myqr}{ \textcolor{black}{\qrcode[height=1.5in]{https://github.com/SergeiSa/Advanced-Control-2026}}
}

\newcommand{\myqrframe}{
	\begin{frame}
		\centerline{Lecture slides are available via Github:}
		\bigskip
		\centerline{\mygit}
		\bigskip
		\myqr
	\end{frame}
}


\newcommand{\bref}[2] {\textcolor{blue}{\href{#1}{#2}}}

%%%%%%%%%%%%%%%%%%%%%%%%%%%%
% code settings

\usepackage{listings}
\usepackage{color}
% \definecolor{mygreen}{rgb}{0,0.6,0}
% \definecolor{mygray}{rgb}{0.5,0.5,0.5}
\definecolor{mymauve}{rgb}{0.58,0,0.82}
\lstset{ 
	backgroundcolor=\color{white},   % choose the background color; you must add \usepackage{color} or \usepackage{xcolor}; should come as last argument
	basicstyle=\footnotesize,        % the size of the fonts that are used for the code
	breakatwhitespace=false,         % sets if automatic breaks should only happen at whitespace
	breaklines=true,                 % sets automatic line breaking
	captionpos=b,                    % sets the caption-position to bottom
	commentstyle=\color{mygreen},    % comment style
	deletekeywords={...},            % if you want to delete keywords from the given language
	escapeinside={\%*}{*)},          % if you want to add LaTeX within your code
	extendedchars=true,              % lets you use non-ASCII characters; for 8-bits encodings only, does not work with UTF-8
	firstnumber=0000,                % start line enumeration with line 0000
	frame=single,	                   % adds a frame around the code
	keepspaces=true,                 % keeps spaces in text, useful for keeping indentation of code (possibly needs columns=flexible)
	keywordstyle=\color{blue},       % keyword style
	language=Octave,                 % the language of the code
	morekeywords={*,...},            % if you want to add more keywords to the set
	numbers=left,                    % where to put the line-numbers; possible values are (none, left, right)
	numbersep=5pt,                   % how far the line-numbers are from the code
	numberstyle=\tiny\color{mygray}, % the style that is used for the line-numbers
	rulecolor=\color{black},         % if not set, the frame-color may be changed on line-breaks within not-black text (e.g. comments (green here))
	showspaces=false,                % show spaces everywhere adding particular underscores; it overrides 'showstringspaces'
	showstringspaces=false,          % underline spaces within strings only
	showtabs=false,                  % show tabs within strings adding particular underscores
	stepnumber=2,                    % the step between two line-numbers. If it's 1, each line will be numbered
	stringstyle=\color{mymauve},     % string literal style
	tabsize=2,	                   % sets default tabsize to 2 spaces
	title=\lstname                   % show the filename of files included with \lstinputlisting; also try caption instead of title
}


%%%%%%%%%%%%%%%%%%%%%%%%%%%%
% URL settings
\hypersetup{
	colorlinks=false,
	linkcolor=blue,
	filecolor=blue,      
	urlcolor=blue,
}

%%%%%%%%%%%%%%%%%%%%%%%%%%

%%%%%%%%%%%%%%%%%%%%%%%%%%%%
% tikz settings

\usepackage{tikz}
\tikzset{every picture/.style={line width=0.75pt}}
\usetikzlibrary{shapes.geometric, decorations.markings, positioning, patterns, calc, arrows.meta}


\subtitle{Advanced Control}
\author{by Sergei Savin}
\centering
\date{\mydate}



\title{Policy Gradient Methods}
\begin{document}
\maketitle




\begin{frame}{Log trick}
	\begin{flushleft}
		
		In this lecture, we will use two properties of the logarithm.
		
		First:
		\[
		\frac{d}{dx}\log f(x)
		=
		\frac{1}{f(x)}\frac{df(x)}{dx}.
		\]
		
		Therefore:
		\[
		\boxed{
			\frac{df(x)}{dx}
			=
			f(x)\frac{d}{dx}\log f(x)
		}
		\]
		
		The same identity applies to gradients with respect to
		a vector of parameters.
		
	\end{flushleft}
\end{frame}


\begin{frame}{Trajectories}
	\begin{flushleft}
		
		A discrete trajectory $\tau$ is a sequence of states
		and actions:
		
		\[
		\tau=(s_0,a_0,s_1,a_1,\ldots,s_{T-1},a_{T-1},s_T).
		\]
		
		In learning, we obtain a trajectory by running an episode:
		
		\begin{itemize}
			\item Sample an initial state $s_0$.
			\item Sample an action $a_t$ according to
			the current policy $\pi_\theta(a_t|s_t)$.
			\item Let the environment sample the next state
			$s_{t+1}$ according to
			$P(s_{t+1}|s_t,a_t)$.
		\end{itemize}
		
		\vspace{0.2cm}
		
		\begin{itemize}
			\item $\pi_\theta(a_t|s_t)$ determines the probability
			of an action given the current state.
			\item $P(s_{t+1}|s_t,a_t)$ determines the probability
			of the next state given the current state and action.
		\end{itemize}
		
	\end{flushleft}
\end{frame}


\begin{frame}{Reward for a trajectory}
	\begin{flushleft}
		
		Knowing a trajectory $\tau$, we can compute the
		total reward associated with it:
		
		\[
		\boxed{
			R(\tau)=\sum_{t=0}^{T-1}r(s_t,a_t)
		}
		\]
		
		\begin{itemize}
			\item This is a deterministic quantity for a fixed
			trajectory.
			\item It is easy to compute from collected data:
			we simply sum the rewards observed during an episode.
		\end{itemize}
		
	\end{flushleft}
\end{frame}


\begin{frame}{Probability of a trajectory}
	\begin{flushleft}
		
		Some trajectories are more likely than others,
		given a specific policy.
		
		We can compute the probability of a trajectory as:
		
		\[
		\boxed{
			p_\theta(\tau)
			=
			p(s_0)
			\prod_{t=0}^{T-1}
			\pi_\theta(a_t|s_t)
			P(s_{t+1}|s_t,a_t)
		}
		\]
		
		where $p(s_0)$ is the probability of the initial state.
		
		Each term describes one step of the trajectory:
		the policy selects an action, and the environment
		samples the next state.
		
	\end{flushleft}
\end{frame}


\begin{frame}{Simplification via logarithm}
	\begin{flushleft}
		
		Remember that the logarithm transforms products into sums:
		
		\[
		\log(fg)=\log f+\log g.
		\]
		
		We can therefore write the log-probability of a trajectory as:
		
		\[
		\boxed{
			\begin{aligned}
				\log p_\theta(\tau)
				={}&\log p(s_0)\\
				&+\sum_{t=0}^{T-1}\log\pi_\theta(a_t|s_t)\\
				&+\sum_{t=0}^{T-1}
				\log P(s_{t+1}|s_t,a_t).
			\end{aligned}
		}
		\]
		
		This expression will make differentiation much easier.
		
	\end{flushleft}
\end{frame}


\begin{frame}{Differentiating the probability of a trajectory}
	\begin{flushleft}
		
		Using the log trick:
		
		\[
		\nabla_\theta p_\theta(\tau)
		=
		p_\theta(\tau)
		\nabla_\theta\log p_\theta(\tau).
		\]
		
		Only the policy terms explicitly depend on $\theta$.
		The initial-state distribution and environment dynamics
		are assumed to be independent of the policy parameters.
		
		Therefore:
		
		\[
		\boxed{
			\nabla_\theta p_\theta(\tau)
			=
			p_\theta(\tau)
			\sum_{t=0}^{T-1}
			\nabla_\theta\log\pi_\theta(a_t|s_t)
		}
		\]
		
	\end{flushleft}
\end{frame}


\begin{frame}{REINFORCE (1/3)}
	\begin{flushleft}
		
		We define the objective as the expected reward:
		
		\[
		\boxed{
			J(\theta)
			=
			\sum_\tau p_\theta(\tau)R(\tau)
		}
		\]
		
		where the summation is over all possible trajectories.
		
		Since the reward of a fixed trajectory does not
		explicitly depend on the policy parameters:
		
		\[
		\nabla_\theta J(\theta)
		=
		\sum_\tau
		\nabla_\theta p_\theta(\tau)R(\tau).
		\]
		
	\end{flushleft}
\end{frame}


\begin{frame}{REINFORCE (2/3)}
	\begin{flushleft}
		
		Substituting our previous identity:
		
		\[
		\begin{aligned}
			\nabla_\theta J(\theta)
			=
			\sum_\tau
			&p_\theta(\tau)R(\tau)\\
			&\cdot
			\sum_{t=0}^{T-1}
			\nabla_\theta\log\pi_\theta(a_t|s_t).
		\end{aligned}
		\]
		
		Equivalently, using the definition of expectation:
		
		\[
		\boxed{
			\nabla_\theta J(\theta)
			=
			\mathbb{E}_{\tau\sim p_\theta}
			\left[
			R(\tau)
			\sum_{t=0}^{T-1}
			\nabla_\theta\log\pi_\theta(a_t|s_t)
			\right]
		}
		\]
		
	\end{flushleft}
\end{frame}


\begin{frame}{REINFORCE (3/3)}
	\begin{flushleft}
		
		We can estimate the gradient using a sampled trajectory:
		
		\[
		\widehat{\nabla_\theta J}
		=
		R(\tau)
		\sum_{t=0}^{T-1}
		\nabla_\theta\log\pi_\theta(a_t|s_t).
		\]
		
		We can then update the policy parameters using
		gradient ascent:
		
		\[
		\boxed{
			\theta
			\leftarrow
			\theta+
			\alpha R(\tau)
			\sum_{t=0}^{T-1}
			\nabla_\theta\log\pi_\theta(a_t|s_t)
		}
		\]
		
		This is the REINFORCE algorithm.
		
	\end{flushleft}
\end{frame}



\begin{frame}{From trajectory reward to reward-to-go}
	\begin{flushleft}
		
		Currently, we assign the same total trajectory reward
		$R(\tau)$ to every action in the trajectory.
		
		Consider an action taken at time $t$.
		Rewards obtained before that action cannot have been
		influenced by it.
		
		Therefore, we can replace the total reward with
		the reward-to-go:
		
		\[
		\boxed{
			G_t=\sum_{k=t}^{T-1}r(s_k,a_k)
		}
		\]
		
		The policy-gradient expression becomes:
		
		\[
		\boxed{
			\nabla_\theta J(\theta)
			=
			\mathbb{E}
			\left[
			\sum_{t=0}^{T-1}
			G_t\nabla_\theta\log\pi_\theta(a_t|s_t)
			\right]
		}
		\]
		
	\end{flushleft}
\end{frame}


\begin{frame}{From reward-to-go to advantage}
	\begin{flushleft}
		
		The state-value and action-value functions are:
		
		\[
		V^\pi(s_t)=\mathbb{E}[G_t|s_t],
		\]
		
		\[
		Q^\pi(s_t,a_t)=\mathbb{E}[G_t|s_t,a_t].
		\]
		
		We can use them to define the advantage:
		
		\[
		\boxed{
			A^\pi(s_t,a_t)
			=
			Q^\pi(s_t,a_t)-V^\pi(s_t)
		}
		\]
		
		The advantage measures how much better an action
		is compared with the average action in that state.
		
		We can replace reward-to-go with an advantage estimate
		without changing the expected policy gradient:
		
		\[
		\boxed{
			\nabla_\theta J(\theta)
			=
			\mathbb{E}
			\left[
			\sum_{t=0}^{T-1}
			A^\pi(s_t,a_t)
			\nabla_\theta\log\pi_\theta(a_t|s_t)
			\right]
		}
		\]
		
	\end{flushleft}
\end{frame}


\begin{frame}{The probability ratio}
	\begin{flushleft}
		
		Suppose we collected data using an old policy
		$\pi_{\theta_{\mathrm{old}}}$, and we are now optimizing
		a new policy $\pi_\theta$.
		
		We define the probability ratio:
		
		\[
		\boxed{
			r_t(\theta)
			=
			\frac{
				\pi_\theta(a_t|s_t)
			}{
				\pi_{\theta_{\mathrm{old}}}(a_t|s_t)
			}
		}
		\]
		
		This ratio tells us how much the new policy changes
		the probability of an action compared with the old policy.
		
		\begin{itemize}
			\item $r_t(\theta)>1$: the new policy makes the action
			more likely.
			\item $r_t(\theta)<1$: the new policy makes the action
			less likely.
			\item $r_t(\theta)=1$: the probability is unchanged.
		\end{itemize}
		
	\end{flushleft}
\end{frame}


\begin{frame}{PPO clipping}
	\begin{flushleft}
		
		We want to prevent excessively large policy changes.
		
		PPO uses the clipped objective:
		
		\[
		\boxed{
			\begin{aligned}
				J^{\mathrm{CLIP}}(\theta)
				=
				\mathbb{E}_{\mathrm{old}}
				\Big[
				\min\big(&r_t(\theta)A_t,\\
				&\operatorname{clip}
				(r_t(\theta),1-\epsilon,1+\epsilon)A_t
				\big)
				\Big].
			\end{aligned}
		}
		\]
		
		\begin{itemize}
			\item The clipping limits the incentive to change
			action probabilities too much.
			\item $\epsilon$ determines the clipping interval.
			\item The objective is maximized with respect to $\theta$.
		\end{itemize}
		
	\end{flushleft}
\end{frame}



\myqrframe






\end{document}
